% TOP_PROBLEM: {"id": 2302027, "title": "Research Problems in Function Theory — Problem 2.27", "category_id": 8, "category": {"id": 8, "name": "analysis", "display_name": "Analysis", "description": "Limits, continuity, calculus, and function theory.", "slug": "analysis", "order_index": 8, "created_at": "2026-07-31T15:26:25.671Z"}, "status": "open", "problem_number": "AMR-022-2027", "set_id": 13}
% FIRST_POSTED: 2026-10-02
% ENTRY_KIND: statement_only
% UnsolvedMath ID 2302027; source catalogue code AMR-022-2027
% !TeX program = lualatex
% Definition and sources only; this file is not an LLM solution attempt.
\documentclass[11pt,a4paper]{article}
\usepackage[margin=25mm]{geometry}
\usepackage{fontspec}
\setmainfont{lmroman10-regular.otf}[BoldFont=lmroman10-bold.otf,
 ItalicFont=lmroman10-italic.otf,BoldItalicFont=lmroman10-bolditalic.otf]
\usepackage{amsmath,amssymb,mathtools}
\usepackage{unicode-math}
\setmathfont{latinmodern-math.otf}
\AtBeginDocument{\let\setminus\smallsetminus}
\usepackage{xurl,hyperref}
\hypersetup{colorlinks=true,urlcolor=blue}
\setcounter{secnumdepth}{0}
\setlength{\parindent}{0pt}
\setlength{\parskip}{5pt}
\setlength{\emergencystretch}{3em}
\begin{document}
\section*{Research Problems in Function Theory — Problem 2.27}\label{2302027}
{\small UnsolvedMath ID 2302027 · AMR-022-2027 · Analysis · AMR Open Problem Lists}
\subsection{Definitions and mathematical statement}
Let $\mathcal E$ be the ring of entire functions of one complex variable. Define a subring $\mathcal A\subset\mathcal E$ by declaring $\phi\in\mathcal A$ when
\[
\phi(z)=\frac{\sum_{j=1}^{p}\exp u_j(z)}{\sum_{k=1}^{q}\exp v_k(z)},
\qquad u_j,v_k\in\mathcal E,\quad p,q<\infty,
\]
where the denominator is not identically zero and all apparent poles are removable. The zero function is allowed; complex constant coefficients can be absorbed into exponential terms, using cancellation when necessary.
Does there exist an entire function $f\notin\mathcal A$ that is integral over $\mathcal A$? In concrete terms, can one find $n\ge2$ and $\phi_1,\ldots,\phi_n\in\mathcal A$ such that
\[
f^n+\phi_1f^{n-1}+\cdots+\phi_n=0
\quad\text{identically on }\mathbb C,
\]
yet $f$ has no quotient-of-finite-exponential-sums representation of the displayed kind?
The source singles out the quadratic case and the case of linear exponent functions as possible starting points. The monic leading coefficient and the requirement that $f$ be entire are both essential. Merely obtaining a multivalued algebraic branch, or allowing poles in $f$, does not provide the requested example.
\subsection{Short English statement}
Is every entire function satisfying a monic algebraic equation over quotients of finite exponential sums itself such a quotient?
\subsection{Sources}
\begin{itemize}
\item[S1] ULAM.AI and the UnsolvedMath Contributors, supplied problems.json, record 2302027 (AMR-022-2027), AMR Open Problem Lists. Catalogue identity and source transcription; CC BY 4.0.
\url{https://huggingface.co/datasets/ulamai/UnsolvedMath}.
\item[S2] Hayman - Research Problems in Function Theory (2018), Problem 2.27. Original source indexed by AMR-022-2027.
\url{https://arxiv.org/abs/1809.07200}.
\end{itemize}
\end{document}
